\documentclass[12pt, a4paper]{article}

% ==========================================
% 1. CODIFICACIÓN Y TIPOGRAFÍA
% ==========================================
\usepackage[utf8]{inputenc}
\usepackage[T1]{fontenc}
\usepackage[spanish, es-tabla]{babel} % 'es-tabla' cambia "Cuadro" por "Tabla" automáticamente
\usepackage{amsmath, amssymb}         % Paquetes matemáticos antes de newtxmath
\usepackage{newtxtext, newtxmath}     % Tipografía Times New Roman (12pt según bases)
\usepackage{microtype}                % Mejora la justificación y estética tipográfica

% ==========================================
% 2. MARGENES E INTERLINEADO (SEGÚN RÚBRICA)
% ==========================================
\usepackage[margin=2.5cm]{geometry}
\usepackage{setspace}
\onehalfspacing                       % Interlineado 1.5 obligatorio
\setlength{\parindent}{0.5cm}         % Sangría de 0.5 cm

% ==========================================
% 3. MATEMÁTICAS, UNIDADES Y TABLAS
% ==========================================
\usepackage{siunitx}                  % Para unidades físicas (ej. \SI{25}{\degree\Celsius})
\sisetup{output-decimal-marker = {,}}  % Coma decimal para español

\usepackage[table]{xcolor}            % Colores con soporte para tablas
\usepackage{booktabs}                 % Tablas profesionales (\toprule, \midrule, \bottomrule)
\usepackage{tabularx}
\usepackage{array}
\usepackage{makecell}
\usepackage{adjustbox}                % Ajuste de tablas anchas
\usepackage{multirow}                 % Necesario para combinar filas en la tabla
\usepackage{ragged2e}

% ==========================================
% 4. FIGURAS Y GRÁFICOS
% ==========================================
\usepackage{graphicx}
\graphicspath{{imagenes/}}
\usepackage{float}
\usepackage{caption}
\usepackage{subcaption}
\captionsetup{font=small, labelfont=bf, justification=centering}

\usepackage{tikz}
\usetikzlibrary{positioning, arrows.meta, babel}

% ==========================================
% 5. BIBLIOGRAFÍA (APA 7) Y ENLACES
% ==========================================
\usepackage{csquotes}
\usepackage[style=apa, backend=biber]{biblatex}
\addbibresource{referencias.bib}

\usepackage[colorlinks=true, linkcolor=black, urlcolor=black, citecolor=black]{hyperref}
\usepackage{bookmark}

% ==========================================
% 6. ENCABEZADOS Y PIE DE PÁGINA
% ==========================================
\usepackage{fancyhdr}
\pagestyle{fancy}
\setlength{\headheight}{15pt}
\fancyhf{}
\fancyhead[L]{\textbf{\coursename} (\coursecode)}
\fancyhead[R]{\small Entrega: \deliverydate}
\fancyfoot[C]{\thepage}
\renewcommand{\headrulewidth}{0.4pt}

% ==========================================
% 7. COMANDOS PERSONALIZADOS
% ==========================================
% Encabezados rotados a 90° para tablas estrechas
\newcommand{\rot}[1]{\rotatebox[origin=l]{90}{\parbox{2.8cm}{\raggedright #1}}}

% ==========================================
% 8. DATOS DEL DOCUMENTO
% ==========================================
\newcommand{\univname}{Pontificia Universidad Católica de Chile}
\newcommand{\schoolname}{Escuela de Ingeniería}
\newcommand{\deptname}{Departamento de Ingeniería Mecánica y Metalúrgica}

\newcommand{\coursename}{Taller de Diseño de Productos Mecánicos}
\newcommand{\coursecode}{ICM3390}
\newcommand{\projecttitle}{Desafío 2: Acondicionador térmico de líquidos de sobremesa}

\newcommand{\authorone}{Juan Francisco Ruz Aedo}
\newcommand{\authortwo}{Francisco Joaquín Retamal Suárez}
\newcommand{\authorthree}{Nicolás Rodrigo Meyer Sironvalle}
\newcommand{\authorfour}{Adolfo Claudio Leiva Mendoza}

\newcommand{\deliverydate}{26 de octubre de 2026}

% ==========================================
% INICIO DEL DOCUMENTO
% ==========================================
\begin{document}

% ------------------------------------------
% PORTADA
% ------------------------------------------
\begin{titlepage}
    \centering
    
    % Logo de la Universidad (Asegúrate de tener 'logo-uc.jpg' en tu carpeta o carpeta imagenes/)
    \includegraphics[width=0.2\textwidth]{logo-uc.jpg}\par
    \vspace{0.5cm}
    
    {\Large \scshape \univname \par}
    \vspace{0.2cm}
    {\large \scshape \schoolname \par}
    \vspace{0.1cm}
    {\normalsize \scshape \deptname \par}
    
    \vspace{0.5cm}
    \rule{\linewidth}{1pt} \par
    \vspace{0.5cm}

    {\Huge \bfseries \projecttitle \par}
    \vspace{0.3cm}
    {\Large Grupo 6 \par}
    \vspace{0.3cm}
    {\large \coursename{} -- \coursecode \par}
    
    \vspace{0.5cm}
    \rule{\linewidth}{1pt} \par

    \vfill 
    
    % Distribución de Integrantes y Docentes
    \begin{minipage}[t]{0.5\textwidth}
        \begin{flushleft} \large
            \textbf{Integrantes:} \\
            \authorone \\
            \authortwo \\
            \authorthree \\
            \authorfour
        \end{flushleft}
    \end{minipage}
    \hfill
    \begin{minipage}[t]{0.4\textwidth}
        \begin{flushright} \large
            \textbf{Profesores:} \\
            Juan Pablo Oyarzun Alcayaga \\
            Ariel Orlando Espinoza Jara
        \end{flushright}
    \end{minipage}

    \vspace{1.5cm}

    {\large \deliverydate \par}
\end{titlepage}

% --- CONFIGURACIÓN DE NUMERACIÓN DESDE EL ÍNDICE ---
\pagenumbering{arabic}
\setcounter{page}{1}

% ------------------------------------------
% ÍNDICES
% ------------------------------------------
\tableofcontents
\newpage
\listoffigures
\listoftables
\newpage

% ==============================================================================
% SECCIONES DEL INFORME (Estructura oficial exigida)
% ==============================================================================

\section{Introducción y contexto de uso}

\section{Necesidades, especificaciones y selección del concepto}
\subsection{Usuario}
Para dar inicio a la elaboración del desafío, se define como usuario a aquellas personas que requieran enfriar bebestibles individuales o calentar alimentos de pequeñas porciones en su hogar, oficina o lugar de estudio. Para esto, se realizan entrevistas a 12 \textit{(modificar el número final cuando se tenga)} personas con el objetivo de identificar las necesidades y preferencias comunes de los usuarios.

\subsection{Necesidades}
Finalizadas las entrevistas, se definen las necesidades del usuario:

\begin{enumerate}
    \item \textbf{Enfriamiento efectivo:} Alcanzar la temperatura ideal de consumo para latas de bebida.
    \item \textbf{Calentamiento efectivo:} Elevar el agua a la temperatura adecuada para la sopa.
    \item \textbf{Tiempo de respuesta:} Completar el acondicionamiento térmico en un rango de tiempo breve.
    \item \textbf{Facilidad de uso:} Operación simple e intuitiva para el usuario.
    \item \textbf{Limpieza e higiene:} Limpieza rápida y accesible de las piezas en contacto con alimentos.
    \item \textbf{Portabilidad:} Diseño compacto y transportable para entornos domésticos o de oficina.
    \item \textbf{Inocuidad:} Agua apta y segura para consumo (sin alteración de sabor u olor).
    \item \textbf{Protección integral:} Aislamiento seguro contra riesgos térmicos, mecánicos y eléctricos.
    \item \textbf{Bajo costo:} Presupuesto de desarrollo y fabricación accesible.
    \item \textbf{Eficiencia energética:} Bajo consumo eléctrico durante los ciclos de acondicionamiento.
\end{enumerate}

Una vez identificadas las necesidades, se les solicita a los entrevistados que distribuyan un total de 100 puntos entre las necesidades mencionadas en orden de importancia para ellos. La tabla \ref{tab:necesidades} muestra el promedio de los resultados obtenidos.

\begin{table}[H]
    \centering
    \begin{tabular}{lc}
        \hline
        \textbf{Necesidad} & \textbf{Puntos} \\
        \hline
        Enfriamiento efectivo & - \\
        Calentamiento efectivo & - \\
        Tiempo de respuesta & - \\
        Facilidad de uso & - \\
        Limpieza e higiene & - \\
        Portabilidad & - \\
        Inocuidad & - \\
        Protección integral & - \\
        Bajo costo & - \\
        Eficiencia energética & - \\
        \hline
    \end{tabular}
    \caption{Promedio de puntos asignados a cada necesidad por los entrevistados.}
    \label{tab:necesidades}
\end{table}

\subsection{Competencia}
En base a las entrevistas realizadas, se identifican productos que son catalogados como competencia. Posterior a una investigación de los productos existentes en el mercado, se realiza una comparación de desempeño en la satisfacción de las necesidades de los usuarios.

\subsection{Especificaciones técnicas}
Las necesidades mencionadas se traducen en especificaciones técnicas. Estas buscan asegurar el cumplimiento de las necesidades de los usuarios mediante objetivos, criterios y métricas, los cuales permiten medir el desempeño del producto en cada una de ellas. Esto significa que para cada especificación existe un objetivo deseado, un criterio de éxito y una métrica de evaluación. Para esto, se agrupan en diferentes categorías para luego exponer una explicación breve y concisa de cada una de ellas.

\begin{itemize}
    \item Desempeño térmico
    \begin{enumerate}
        \item Enfriamiento efectivo: .
        \item Calentamiento efectivo: .
        \item Tiempo de respuesta: .
    \end{enumerate}
    \item Uso e interacción
    \begin{enumerate}
        \item Facilidad de uso: .
        \item Limpieza e higiene: .
        \item Portabilidad: .
    \end{enumerate}
    \item Seguridad
    \begin{enumerate}
        \item Inocuidad: .
        \item Protección integral al usuario: .
    \end{enumerate}
    \item Viabilidad y costo
    \begin{enumerate}
        \item Bajo costo: .
        \item Bajo consumo energético: .
    \end{enumerate}
\end{itemize}

En base a las justificaciones anteriores se definen las especificaciones técnicas, con sus criterios, objetivos y métricas asociados en la Tabla \ref{tab:especificaciones}.

\begin{table}[htbp]
\centering
\small
\setlength{\tabcolsep}{6pt}
\begin{tabularx}{\textwidth}{@{}>{\RaggedRight\arraybackslash}X >{\Centering\arraybackslash}p{2.5cm} >{\Centering\arraybackslash}p{1.5cm} >{\Centering\arraybackslash}p{1.0cm}@{}}
\toprule
\textbf{Especificación Técnica} & \textbf{Objetivo} & \textbf{Unidad} & \textbf{Sentido} \\ 
\midrule

% --- 1. ENFRIAMIENTO EFECTIVO ---
\multicolumn{4}{l}{\textbf{Necesidad:} 1. Enfriamiento efectivo} \\
\multicolumn{4}{l}{\textit{Criterio:} Desempeño Térmico (Modo Frío)} \\
\nopagebreak
Temperatura objetivo de la lata ($T_{f,c}$) & $\le 8$ & \si{\celsius} & $\downarrow$ \\
Volumen nominal de la bebida ($V_c$) & $354$ & \si{\ml} & $\uparrow$ \\
\midrule

% --- 2. CALENTAMIENTO EFECTIVO ---
\multicolumn{4}{l}{\textbf{Necesidad:} 2. Calentamiento efectivo} \\
\multicolumn{4}{l}{\textit{Criterio:} Desempeño Térmico (Modo Calor)} \\
\nopagebreak
Temperatura objetivo del agua ($T_{f,h}$) & $\ge 85$ & \si{\celsius} & $\uparrow$ \\
Volumen de agua para sopa ($V_h$) & $300$ & \si{\ml} & $\uparrow$ \\
\midrule

% --- 3. TIEMPO DE RESPUESTA ---
\multicolumn{4}{l}{\textbf{Necesidad:} 3. Tiempo de respuesta} \\
\multicolumn{4}{l}{\textit{Criterio:} Cinética de Transferencia Térmica} \\
\nopagebreak
Tiempo de enfriamiento ($t_c$) & $\le 15$ & \si{\minute} & $\downarrow$ \\
Tiempo de calentamiento ($t_h$) & $\le 10$ & \si{\minute} & $\downarrow$ \\
\midrule

% --- 4. FACILIDAD DE USO ---
\multicolumn{4}{l}{\textbf{Necesidad:} 4. Facilidad de uso} \\
\multicolumn{4}{l}{\textit{Criterio:} Usabilidad e Interfaz} \\
\nopagebreak
Pasos requeridos para iniciar el proceso & $\le 2$ & Acciones & $\downarrow$ \\
\midrule

% --- 5. LIMPIEZA E HIGIENE ---
\multicolumn{4}{l}{\textbf{Necesidad:} 5. Limpieza e higiene} \\
\multicolumn{4}{l}{\textit{Criterio:} Mantenibilidad y Sanidad} \\
\nopagebreak
Tiempo de desmontaje y lavado de piezas & $\le 3$ & \si{\minute} & $\downarrow$ \\
\midrule

% --- 6. PORTABILIDAD ---
\multicolumn{4}{l}{\textbf{Necesidad:} 6. Portabilidad} \\
\multicolumn{4}{l}{\textit{Criterio:} Ergonomía y Dimensiones} \\
\nopagebreak
Masa total del dispositivo ($m$) & $\le 4,0$ & \si{\kilogram} & $\downarrow$ \\
Envolvente máxima ($A \times L \times H$) & $\le 30 \times 30 \times 35$ & \si{\cm} & $\downarrow$ \\
\midrule

% --- 7. INOCUIDAD ---
\multicolumn{4}{l}{\textbf{Necesidad:} 7. Inocuidad} \\
\multicolumn{4}{l}{\textit{Criterio:} Compatibilidad de Materiales} \\
\nopagebreak
Compatibilidad alimentaria de superficies en contacto con agua & Grado Alimentario & Cualitativo & $\rightarrow$ \\
\midrule

% --- 8. PROTECCIÓN INTEGRAL ---
\multicolumn{4}{l}{\textbf{Necesidad:} 8. Protección integral} \\
\multicolumn{4}{l}{\textit{Criterio:} Seguridad Térmica, Mecánica y Eléctrica} \\
\nopagebreak
Temperatura máxima externa de carcasa ($T_{\text{ext}}$) & $\le 40$ & \si{\celsius} & $\downarrow$ \\
Protecciones eléctricas y contra humedad & IPX2 / Aislamiento & Grado IP & $\uparrow$ \\
\midrule

% --- 9. BAJO COSTO ---
\multicolumn{4}{l}{\textbf{Necesidad:} 9. Bajo costo} \\
\multicolumn{4}{l}{\textit{Criterio:} Factibilidad Económica} \\
\nopagebreak
Costo total de desarrollo y materiales (BOM) & $\le 50.000$ & CLP & $\downarrow$ \\
\midrule

% --- 10. EFICIENCIA ENERGÉTICA ---
\multicolumn{4}{l}{\textbf{Necesidad:} 10. Eficiencia energética} \\
\multicolumn{4}{l}{\textit{Criterio:} Gestión Energética} \\
\nopagebreak
Potencia eléctrica máxima consumida ($P_{\text{máx}}$) & $\le 120$ & \si{\watt} & $\downarrow$ \\

\bottomrule
\end{tabularx}
\caption{Matriz de especificaciones técnicas de ingeniería estructurada según las necesidades del usuario.}
\label{tab:especificaciones}
\end{table}



\subsection{Selección del concepto}


\section{Arquitectura y modelos desarrollados}

\section{Simulaciones y justificación físico-matemática}

\section{Caracterización, calibración y validación inicial}

\section{Desarrollo del producto final y plan de fabricación}

\section{Resultados y discusión}

\section{Conclusiones}

% ------------------------------------------
% BIBLIOGRAFÍA
% ------------------------------------------
\newpage
\nocite{*}
\printbibliography[heading=bibintoc, title={Referencias Bibliográficas}]

\end{document}